\scanpage{208}
\mspage{124}
% archive sheet number: 207
% element: text | top, lines 1-3 | C'est essentiellement
\unsure{C'est|(C'est} essentiellement la \unsure{même} démonstration, légèrement plus simple. On aurait dû commencer par là.)

% element: text | upper third, line 4 | Corollaire 4.5.12
\uline{Corollaire} 4.5.12. Si $\pi'\geq\pi\geq\pi_1$, on a \inserted{pour $x\in\mathrm{Ob}\,M_{\pi'}$}
% element: display | upper third, line 5 | cardinal slice skeleton
\[
\operatorname{card}(M_\pi/x)_0\leq\pi'{}^{(\pi^{\pi_0})}
\tag*{(4.5.12.1)}
\]
% element: text | upper third, line 6 | En effet
En effet,
% element: display | upper middle, lines 6-9 | annotated object count
\[
\operatorname{card}\mathrm{Ob}(M_\pi/x)_0
\underset{\text{par 4.5.11}}{\leq}
\underbrace{
\underbrace{\operatorname{card}\mathrm{Ob}(M_\pi)_0}_{\substack{\leq2^{\pi^{\pi_0}}\ \text{par}\\4.5.10}}
\times(\pi'{}^{\pi_0})^{\pi^{\pi_0}}
}_{(\deleted{\unsure{2^{\pi^{\pi_0}}|2\pi}}=\pi'{}^{(\pi^{\pi_0})})}.
\]
% element: text | middle, line 10 | d'où cardinal Ob
d'où
% element: display | middle, line 10 | resulting object bound
\[
\operatorname{card}\mathrm{Ob}(M_\pi/x)_0\leq\pi'{}^{(\pi^{\pi_0})},
\]
% element: text | middle, line 11 | d'où cardinal Fl
d'où
% element: display | middle, line 11 | arrow count bound
\[
\operatorname{card}\mathrm{Fl}(M_\pi/x)_0
\leq\pi'{}^{(\pi^{\pi_0})}\times\pi'{}^{(\pi^{\pi_0})}
\times(\pi^{\pi_0})^{\pi^{\pi_0}}
\leq\pi'{}^{(\pi^{\pi_0})},
\]
% element: text | middle, line 12 | Cela donne donc
Cela donne donc (\unsure{reven. à|revenant à|revenons à} : 4.5.9.1)
% element: display | lower middle, line 13 | cardinal I zero
\[
\operatorname{card}\unsure{I|\bar I}_0\leq\pi'{}^{(\pi^{\pi_0})},
\]
% element: text | lower third, lines 14-18 | D'autre part
\deleted{\unsure{et}} D'autre part, $\unsure{I'_0|I'}$ peut se définir en termes de $I_0$ comme l'ens des sous-catégories de $I_0$ qui sont \deleted{grandes devant $\pi_0$} de cardinal $\leq\pi$, donc leur nombre est majoré par
% element: display | lower third, line 19 | power bound for number of subcategories
\[
(\operatorname{card}I_0)^\pi
\leq(\pi'{}^{\pi^{\pi_0}})^\pi
=\pi'{}^{\pi\pi^{\pi_0}}
=\pi'{}^{(\pi^{\pi_0})};\quad\deleted{\unsure{qed.|AQT}}
\]
% element: deletion | foot, crossed heading | Corollaire 4.5.13 Finalement
\deleted{\uline{\unsure{Corollaire}} 4.5.13. Finalement,}
% element: insertion | foot, above crossed heading | Mais ce n'est pas tout
\inserted{Mais ce n'est pas tout. \unsure{pas la|par la|pour la}}
% element: text | foot, lines 21-24 | majoration naïve
majoration naïve, puisque
% element: display | foot, line 23 | comparison with naïve bound
\[
\pi'{}^\pi\leq\unsure{\pi'{}^{(\pi^{\pi_0})}|\pi^{(\pi^{\pi_0})}}
\]
% element: text | foot, line 24 | et qu'il est possible
et qu'il est possible, \unsure{a priori|à priori}, que l'inégalité
% element: margin note | lower right, diagonally | NB on a égalité
% slanted
\margin{\uline{NB} On a égalité si $\pi'=2^{c'}$ avec $c'\geq\pi^{\pi_0}$, soit $\pi=2^c$ avec $c\geq\pi_0$.}
